% ===== 原文 pp.39--41 =====
% =====================================================================
%  tr06b.tex —— §6「未来方向」后半 + 致谢（原文 pp.39--41）
%  负责人：tr-06
%  覆盖：原文 p.39 第二段（P14）起，经 p.40、p.41 的致谢，到致谢末句为止。
%        （p.39 的第一段 P13 已由 tr06a.tex 在段落边界处结束，避免分段假象。）
%        不写 References（由 Lead 的 c-refs.tex 负责）。
%  本块无编号公式（§6 全文没有 \tag）、无图、无表；有一个原文脚注（脚注 5）。
% =====================================================================

这在局域化计算中造成了一个很大的困难，因为局域化只适用于理论中具有拉格朗日量
的那部分。幸运的是，我们仍然可以把局域化应用于来自圆柱面的矢量多重态。例如，
仅凭这一点，我们就可以猜想：当 $G$ 是 $A_n$、$D_n$ 或 $E_{n=6,7,8}$ 之一时，
$G$ 规范理论的瞬子配分函数应当由 $W_G$ 代数控制。这在数学界引起了相当不小的
兴趣，如今已被严格证明 [51--53]。然而，$T_G$ 理论的 $S^4$ 配分函数仍属未知。
在 2d/4d 对应关系下，它应当映为 $G$ 型 Toda 理论在一般动量下的三点函数。这
同样未知，除非 $G$ 是 $\SU(2)$ 型。在这里我们看到「困难的守恒定律」在起作用：
对应关系一侧困难的东西，在另一侧也困难。与此同时，我们也可以说：对应关系任何
一侧的突破都会对另一侧产生巨大的影响。作者希望有感兴趣的读者能迈出这样划时代
的一步。$T_G$ 理论的 $S^1\times S^3$ 配分函数仅由结合性就已被推出，至少是在
二维的 $(q,t)$ 切片中。这一信息已被用来更好地理解 $T_G$ 理论。

% [Tutor Note: p.39，这段是 6.3 节的技术核心，逻辑链要看全：(1) 局域化需要
%  拉格朗日量；(2) $G\neq\SU(2)$ 时 $T_G$（三穿刺球面）没有拉格朗日描述；
%  (3) 但圆柱面部分（$G$ 矢量多重态）有拉格朗日量，可以局域化；(4) 于是
%  只能「猜想」$G$ 规范理论的瞬子配分函数由 $W_G$ 代数控制——这正是 AGT
%  对应在一般 $G$ 的版本：瞬子配分函数 = $W_G$ 共形块；(5) 这个猜想已被数学
%  严格证明 [51--53]，但 $T_G$ 本身的 $S^4$ 配分函数（= 一般动量 Toda 三点
%  函数）仍未知。注意 (4) 缺的只是「瞬子部分」：经典因子与单圈因子都能算
%  （任意表示的超多重态，其单圈行列式都有现成公式），真正缺的是三穿刺球面
%  物质的瞬子贡献。对照 §5.2：$G=\SU(2)$
%  时 $W_G$ 就是 Virasoro 代数，这句话退化为原始的 AGT 等式（瞬子配分函数
%  = Liouville 共形块）。]
% [Tutor Note: p.39，「conservation law of the difficulty」是作者对全书处境的
%  一句概括：2d/4d 对应把「算不动的 4d 量」与「算不动的 2d 量」配成一对，
%  所以任何一侧的突破都会立刻给另一侧一个可算的答案；反之只要一侧还难，
%  另一侧也一定还难。本节的例子是「$T_G$ 的 $S^4$ 配分函数难 ↔ 一般动量
%  Toda 三点函数难」；而 $S^1\times S^3$ 指标（不需要 ADHM 式瞬子构造）先
%  在 2015 年有进展，2d 一侧的相应推广也就同步有进展。这句话也可以当成检验
%  「对应关系是否真成立」的启发式：两侧的难度曲线应当同步。]
% [Tutor Note: p.39，「The author hopes that an interested reader will do make
%  such an epoch-making step.」中的 will do make 是原文笔误（已用 700 dpi 裁图
%  复核，原文确为 will do make），应为 will make（或 will do such an
%  epoch-making step）。译文按「迈出这样划时代的一步」处理，意思未变。]

我们可以把一个由同态 $\phi$ 指定的规则余维 2 算符放到 $S^4\times C_2$ 上，
使它缠绕整个 $C_2$，并占据一个 $S^2\subset S^4$。已知这会把 $C_2$
上的 $W_G$ 对称性修改为一个更一般的 W 代数 $W(G,\phi)$，后者由量子
Drinfeld--Sokolov 约化得到。但二维理论本身还没有被做出来。

% [Tutor Note: p.39，$W(G,\phi)$ 与上一段 $W_G$ 的关系：$\phi$ 是
%  $\fsu(2)\to\fg$ 的同态，它给 $\fg$ 一个 $\fsu(2)$ 分级；量子
%  Drinfeld--Sokolov 约化就是按这个分级把仿射代数 $\hat\fg$ 约化，得到 W 代数
%  $W(G,\phi)$。取 $\phi$ 为平凡同态（即前面的「全穿刺点」）时 $W(G,0)=W_G$，
%  正是上一段的 $W_G$ 代数；所以这条 bullet 是「把全穿刺点换成任意规则穿刺点」
%  之后 2d 对称性代数的相应替换。作者说 2d 理论本身「has not been worked
%  out」，意思是：只知道对称性代数是 $W(G,\phi)$，但对应的 2d 作用量/关联函数
%  （Toda 型理论的适当推广）还没写出来——这是 6.3 节留下的另一个具体缺口。
%  另外，本页的 $\phi$ 与下段的 $\varphi$ 在原文中是两个不同字形：前者是同态
%  $\phi$，后者是不变张量 $\varphi:R_1\otimes R_2\otimes R_3\to\C$；高倍裁图
%  可确认（$\phi$ 为带斜直线的椭圆，$\varphi$ 为带尾的卷曲形），不要混为一谈。]

现在让我们简短地讨论 $G$ 型六维理论的余维 4 算符。它们由 $G$ 的不可约表示
标记。当三个表示 $R_{1,2,3}$ 之间存在一个不变张量
$\varphi: R_1\otimes R_2\otimes R_3\to\C$ 时，我们可以考虑沿一条一维轨迹的
三个余维 4 算符的交汇（junction）。我们可以用许多方式把这样一个交汇放进
$X_4\times C_2$；从 $X_4$ 上的四维理论与 $C_2$ 上的二维理论这两个观点来分析
该设定，会产生各种对应关系，其中有些已经被做出来了。

% [Tutor Note: p.39，最后一段的「不变张量 $\varphi$」刻画的是三个余维 4
%  算符沿一条一维轨迹交汇时的耦合：只有当 $R_1\otimes R_2\otimes R_3$ 含
%  单态时，三者才能在交汇处写下一个不变量。最容易想当然的地方是：$\SU(2)$
%  的三个二重态并不可行——$\mathbf{2}^{\otimes3}=\mathbf{4}\oplus
%  \mathbf{2}\oplus\mathbf{2}$ 不含单态，所以「三个二重态标记的余维 4 算符
%  交汇」是写不出不变耦合的；能这样交汇的最简单例子是三个伴随（三重态）
%  $\mathbf{3}$，其不变张量就是结构常数 $\epsilon_{ijk}$。
%  这段的用途：同一个 6d 交汇放在 $X_4\times C_2$ 上，从 $X_4$ 一侧读是 4d
%  算符、从 $C_2$ 一侧读是 2d 共形块，两个观点各给出一种「对应关系」；作者
%  说「some of which have been worked out」即指这类已被算出的情形。
%  另外，一般 $G$ 的 $T_G$ 理论没有拉格朗日描述，是因为它是六维理论在带三个
%  全穿刺点的球面上的紧致化，没有可用的「圆柱面矢量多重态」分解；只有
%  $G=\SU(2)$ 时 $T_2$ 才是自由的三基表示半超多重态理论。]

\subsection{其它时空、其它理论}

在本综述中，我们只讨论了把六维 $\cN{=}(2,0)$ 理论放到形如 $X_4\times C_2$
的一个特定时空上——其中我们沿 $C_2$ 做部分拓扑扭转，以便在 $X_4$ 上得到
$\cN{=}2$ 超对称。事实上，本综述实际用到的选择还要更受限：我们只考虑了两种
情形，$X_4=S^1\times S^3$ 与 $X_4=S^4$。显然，这份很短的清单可以扩充。关于
$X_4=S^1\times S^3/\Z_k$ 与 $X_4=S^4/\Z_k$ 已有一些工作。例如，已知
$\SU(2)$ 型六维 $\cN{=}(2,0)$ 理论放在 $S^4/\Z_2$ 上会导致二维一侧的
$\cN{=}(1,1)$ 超 Liouville 理论。另一个选择是取紧致的复 toric 曲面作为
$X_4$。为此，局域化计算的最终公式已由 Nekrasov 在 [54] 中宣布；其细节最近%
% [原文 p.40]
由不同的作者在 [55] 中给出。不过，以 $X_4$ 为标签的相应二维理论还没有被
理解。

% [Tutor Note: p.39--40，「沿 $C_2$ 做部分拓扑扭转」是全文一切构造的前提：
%  用 $C_2$ 的 $\SO(2)$ 旋转与 R 对称性中的某个 $U(1)$ 做混合（一种「半扭转」），
%  从而对任意 Riemann 曲面 $C_2$ 都能保住 $X_4$ 上的 $\cN{=}2$ 超对称（这就是为什么
%  §3、§4 可以放心地把 $C_2$ 当「内部空间」、在 $X_4$ 上做局域化）。本节
%  列的其实是「$X_4$ 还能换成什么」；括号里的 $S^4/\Z_2$ 例子说明对应关系
%  对 $X_4$ 的几何很敏感：同一个六维理论放在 $S^4$ 上给 Liouville，放在
%  $S^4/\Z_2$ 上就给出 $\cN{=}(1,1)$ 超 Liouville（保留的超对称更少）。]
% [Tutor Note: p.40，toric 曲面这段是「换 $X_4$」的第一个具体进展与第一个
%  具体缺口：Nekrasov 在 [54] 宣布了局域化计算的最终公式——对紧致复 toric
%  曲面，等变局域化归结为对 toric 图各固定点的求和；「细节」[55] 则是把它
%  写成显式公式的一批工作。缺口在于：公式只给出 4d 配分函数这个数，要把它
%  认成「以 $X_4$ 为标签的 2d 理论」，得先知道那个 2d 理论的作用量与对称性；
%  $X_4=S^4$ 时 2d 侧是 Liouville，而对一般 toric $X_4$ 连候选的 2d 理论都
%  还没有——这就是「were not understood yet」的含义。]

我们也可以把非紧的复曲面取作 $X_4$。此时我们得不到一个完全意义上的二维
理论；相反，其上的超对称局域化给出二维手征代数。$X_4=\R^4$ 的情形最初由
[5] 研究，给出的正是我们在 §4.3 中讨论过的瞬子贡献。当 $X_4$ 是 ALE 空间时
也有各种研究，这一方向由 [56] 开创。参考文献的一个样本见例如 [57--76]。%
\footnote{作者感谢 F. Sala 在整理这份参考文献清单时的帮助。}

% [Tutor Note: p.40，脚注：原文此脚注编号为 5（全文脚注从 1 起编号）。本译本的
%  脚注由 LaTeX 自动编号，只要上游各块按原文顺序各就各位，序号应当仍是 5；
%  若上游 chunk 有脚注增删，则序号可能不同，以原文为准。]
% [Tutor Note: p.40，非紧 $X_4$ 这一段有一个重要的概念转折：$X_4$ 非紧时没有
%  可归一化的 4d 涨落，所以得不到「完整的」2d 理论，只剩下一个二维手征代数
%  （chiral algebra，即顶点算子代数 VOA）。$X_4=\R^4$ 就是 §4.3 的
%  $\Omega$-背景与 Nekrasov 瞬子配分函数；把它看成某个 2d 手征代数的配分
%  函数，是后来所谓 4d/2d 手征代数对应的出发点（本条为译者补充，原文只说到
%  「gives rise to 2d chiral algebras」为止）。[56] 开创的 ALE 空间情形同理：
%  非紧 ALE 给出的是手征代数而非 2d 理论。]

当然，我们也可以考虑六维时空的其它分解方式。例如，我们可以考虑把六维
$\cN{=}(2,0)$ 理论放到形如 $X_{6-d}\times C_d$ 的时空上，沿 $C_d$ 做部分
拓扑扭转，使得保留下来的超对称的平方在 $X_{6-d}$ 中生成一个等距变换。注意
这个设定在交换 $d\leftrightarrow 6-d$ 下并不对称。在本综述中我们只讨论了
$d=2$ 的情形；其它情形 $d=3,4$ 已被分析过，主要是针对 $\SU(N)$ 型的六维
$\cN{=}(2,0)$ 理论。例如，读者可以在文集中的 [77] 一文里找到关于 3d-3d
情形的一些讨论。如果把各种余维 2 与余维 4 算符也纳入这个一般设定，我们还能
获得更多乐趣。

% [Tutor Note: p.40，「$d\leftrightarrow 6-d$ 下不对称」的原因：拓扑扭转只沿
%  $C_d$ 做，$X_{6-d}$ 是「物理时空」。交换 $d$ 与 $6-d$ 会把「被扭转的因子」
%  与「物理因子」互换（$d=2$ 时是 $X_4\times C_2$，$d=4$ 时是 $X_2\times C_4$），
%  而扭转只对一个因子做，所以这不是对称操作。顺带说：$d=3$ 就是 3d-3d 对应
%  （六维理论在 3-流形上紧致化 ↔ 3d 理论），$d=2$ 就是本文的 2d/4d 对应；
%  两者是同一思想在不同维数上的实例。]

仅仅把六维 $\cN{=}(2,0)$ 理论放到乘积流形上，就给我们带来了如此大量的、正在
被逐步解决的低维量子场论对应关系。那么一个自然的方向就是：寻找六维
$\cN{=}(2,0)$ 理论之外的其它理论作为出发点。在六维，还有许多其它被认为紫外
完备的 $\cN{=}(1,0)$ 超对称理论。它们比 $\cN{=}(2,0)$ 理论还要不被理解，但
最近已有一些系统地探索其性质的活动。一旦基本特征被理解，把它们放到乘积流形
上应当并非不可能，而这有望带来全新类型的 4d-2d 及其它对应关系。

当然，我们也可以从更低维的、更熟悉的规范理论出发，把它们放到形如
$X_{D-d}\times C_d$ 的乘积流形上。例如，规范群为 $G$ 的 $\cN{=}4$
超杨--米尔斯理论曾被放到形如 $X_2\times C_2$ 的时空上考虑，沿 $C_2$ 做部分
拓扑扭转；这正是 Witten 处理几何 Langlands 对应的出发点。在 $X_2$ 上得到的
二维理论是 $C_2$ 上 $G$-Hitchin 系统模空间上的非线性 sigma 模型；当把 $C_2$
切成两半时，它似乎没有像四维 class S 理论那样漂亮的切割--粘接描述。一个
已知拉格朗日量的四维规范理论，在 $C_2$ 上紧致化后的行为，居然比一个没有
拉格朗日量的神秘六维理论的行为还要复杂，这听起来很奇怪。主要差别似乎在于：
Riemann 曲面 $C_2$ 上规范场可能的超对称构型相当丰富，因为 $C_2$ 可以有各种
非平凡的一维闭链，规范场可以绕这些闭链具有非平凡的和乐。而在六维
$\cN{=}(2,0)$ 理论的情形，我们没有有用的拉格朗日描述；非常朴素地说，人们说
它是由非阿贝尔两形式场构成的理论。无论这句话是什么意思，它都暗示：Riemann
曲面 $C_2$ 上不可能有太多的超对称构型，%
% [原文 p.41]
因为 $C_2$ 中只有一个二维闭链。话虽如此，既然出发点有显式的拉格朗日描述，
原则上就应当能够把 $X_{D-d}$ 上与 $C_d$ 上的低维场论对应关系的方方面面都算
出来——作者认为这值得去做。

% [Tutor Note: p.40--41，「非阿贝尔两形式场」与「只有一个二维闭链」这段是
%  作者的启发式论证，值得补清：超对称构型的模由规范场沿闭链的和乐（Wilson
%  线）等给出；$C_2$ 是复一维，$H_2(C_2)$ 只有一维（由基本类 $[C_2]$ 生成），
%  所以两形式场能绕出的独立模很少。而 4d 规范场在 $C_2$ 上可以绕各种非平凡
%  1-闭链取非平凡和乐（Hitchin 系统 Higgs 场的数据正住在这里），模空间因此
%  大得多：$X_2$ 上的 2d 理论是 Hitchin 模空间上的非线性 sigma 模型，也就
%  没有好的切割--粘接描述。作者用这个对比解释「有拉格朗日量的 4d 规范理论
%  反而更难」这一表面悖论；注意这只是 very naive 的直观，作者自己用了
%  「Whatever this statement means」的谨慎措辞——(2,0) 理论至今没有严格的
%  自由度描述。]

% =====================================================================
% 致谢（用 \section* + \addcontentsline，不要用 \section{致谢}，否则会多编出第 7 节）
% =====================================================================
\section*{致谢}
\addcontentsline{toc}{section}{致谢}

感谢 Tatsuma Nishioka 仔细阅读草稿并提出各种有用的意见。本工作部分受到
JSPS Grant-in-Aid for Scientific Research No.\ 25870159 的资助，部分受到
WPI Initiative, MEXT, Japan（在东京大学 IPMU）的资助。
